%%%%%%%%%%%%%%%%%%%%%%%%%%%%%%%%%%%%%%%%%%%%%%%%%%%%%%%%%%%%%%%%%%%%%%%%%%%%%%%%
%% 
\cleardoubleoddpage%  Make sure to start each chapter on a new odd page
\chapter{Experimental Setup}
\label{sec:setup}

This chapter provides comprehensive technical specifications and implementation details for the experimental evaluation. While Chapter~\ref{sec:methodology} described the conceptual framework and scientific rationale, this chapter focuses on the specific configurations, parameter values, and computational infrastructure used to execute the experiments.

\section{Implementation Environment}
\label{sec:implementation_environment}

\subsection{Software Stack}
\label{sec:software_stack}

The experimental implementation leverages modern deep learning frameworks and scientific computing libraries:

\textbf{Core Deep Learning Framework:}
\begin{itemize}
    \item \textbf{PyTorch 2.x}: Primary deep learning framework providing automatic differentiation and GPU acceleration
    \item \textbf{PyTorch Lightning 2.x}: High-level framework for training organization, distributed computing, and logging infrastructure
    \item \textbf{CUDA 11.8+}: NVIDIA CUDA toolkit for GPU computation
\end{itemize}

\textbf{xLSTM Integration:}
\begin{itemize}
    \item \textbf{xlstm package}: Official implementation providing optimized \ac{sLSTM} and \ac{mLSTM} components with CUDA backend support (\url{https://github.com/NX-AI/xlstm})
\end{itemize}

\textbf{Signal Processing and Data Handling:}
\begin{itemize}
    \item \textbf{NumPy 1.24+}: Numerical computing and array operations
    \item \textbf{SciPy 1.10+}: Signal processing for resampling and filtering
    \item \textbf{MNE-Python 1.3+}: \ac{EEG} data loading and \ac{EDF} format handling
    \item \textbf{PyHealth}: Medical data analysis utilities and evaluation metrics
\end{itemize}

\textbf{Experiment Tracking and Visualization:}
\begin{itemize}
    \item \textbf{Weights \& Biases (wandb)}: Comprehensive experiment tracking, metric logging, and model artifact management
    \item \textbf{Matplotlib/Seaborn}: Visualization for results analysis
\end{itemize}

\subsection{Hardware Configuration}
\label{sec:hardware_config}

The experimental evaluation is conducted on a single-GPU workstation:

\textbf{GPU Specifications:}
\begin{itemize}
    \item \textbf{Model}: NVIDIA GeForce RTX 3090
    \item \textbf{Memory}: 24GB GDDR6X
    \item \textbf{CUDA Cores}: 10,496
    \item \textbf{Tensor Cores}: 328 (3rd generation)
    \item \textbf{Memory Bandwidth}: 936 GB/s
\end{itemize}

The 24GB memory capacity enables processing of long \ac{EEG} sequences and supports the memory-intensive operations required by matrix \ac{LSTM} implementations. The substantial memory bandwidth facilitates efficient gradient computation during backpropagation through long sequences.

\textbf{Training Strategy:} Single-GPU training without distributed computing. While PyTorch Lightning supports multi-GPU distributed training, the dataset size and model complexity make single-GPU training sufficiently efficient for this experimental evaluation.


\section{Model Configuration}
\label{sec:model_configuration}

\subsection{BIOT Tokenization Parameters}
\label{sec:tokenization_parameters}

The biosignal tokenization module employs \ac{STFT} to convert temporal signals into frequency-domain token representations. The \ac{STFT} configuration determines the frequency and temporal resolution of the resulting tokens.

\textbf{STFT Configuration:}
\begin{itemize}
    \item \textbf{FFT Size (\texttt{n\_fft})}: 200 samples
    \item \textbf{Hop Length (\texttt{hop\_length})}: 100 samples
    \item \textbf{Window Function}: Hann window
    \item \textbf{Sampling Rate}: 200 Hz (after resampling)
\end{itemize}

\textbf{Resolution Calculations:}

Frequency resolution:
\begin{equation}
\Delta f = \frac{f_s}{N} = \frac{200 \text{ Hz}}{200} = 1 \text{ Hz}
\end{equation}

Temporal resolution (hop time):
\begin{equation}
\Delta t = \frac{H}{f_s} = \frac{100}{200 \text{ Hz}} = 0.5 \text{ seconds}
\end{equation}

Frame overlap:
\begin{equation}
\text{Overlap} = \frac{N - H}{N} = \frac{200 - 100}{200} = 50\%
\end{equation}

This configuration provides 1 Hz frequency resolution spanning 0-100 Hz (Nyquist limit) with 0.5-second temporal resolution and 50\% overlap between consecutive frames. The 50\% overlap ensures smooth temporal continuity while maintaining computational efficiency.

\textbf{Token Sequence Length:} The tokenization process transforms raw temporal samples into frequency-domain tokens through \ac{STFT}. For a signal window of duration $T$ seconds sampled at 200 Hz, the input contains $L = T \times 200$ time-domain samples. The \ac{STFT} with \texttt{center=False} generates:
\begin{equation}
N_{tokens} = \left\lfloor \frac{L - N}{H} \right\rfloor + 1 = \left\lfloor \frac{T \cdot 200 - 200}{100} \right\rfloor + 1 = 2T - 1
\end{equation}

\textbf{Concrete Examples:}
\begin{itemize}
    \item \textbf{2-2 configuration (5s total):} 1000 samples $\rightarrow$ 9 tokens/channel $\rightarrow$ 144 total tokens
    \item \textbf{3-3 configuration (7s total):} 1400 samples $\rightarrow$ 13 tokens/channel $\rightarrow$ 208 total tokens
    \item \textbf{4-4 configuration (9s total):} 1800 samples $\rightarrow$ 17 tokens/channel $\rightarrow$ 272 total tokens
\end{itemize}

Each token represents a 0.5-second window of frequency content across the 0-100 Hz range, providing a compact sequence representation suitable for transformer processing.

\subsection{Baseline Transformer Configuration}
\label{sec:baseline_transformer_config}

The baseline \ac{BIOT} architecture employs a Linear Attention Transformer as the sequence encoder, serving as the reference configuration against which \ac{xLSTM} variants are compared~\cite{yang2023biot}.

\textbf{Linear Attention Mechanism:} The baseline uses linear attention rather than standard softmax attention, reducing computational complexity from quadratic to linear in sequence length. This enables efficient processing of longer \ac{EEG} sequences while maintaining the Transformer's ability to model long-range dependencies. The linear attention approximation computes attention scores without the expensive softmax operation, trading some representational capacity for computational efficiency.

\textbf{Architecture Parameters:}
\begin{itemize}
    \item \textbf{Embedding Dimension}: 256 (matching \ac{xLSTM} variants for fair comparison)
    \item \textbf{Num Blocks}: 8 (stacked transformer layers)
    \item \textbf{Attention Heads}: 8 (for multi-head attention mechanism)
    \item \textbf{Context Length}: 1024 tokens (maximum sequence length)
    \item \textbf{Dropout}: 0.2
\end{itemize}

\textbf{Layer Composition:} Each transformer block consists of a linear multi-head attention layer followed by a position-wise feedforward network, with residual connections and layer normalization applied around both sub-layers. This follows the standard Transformer architecture pattern but substitutes linear attention for computational efficiency.

For complete implementation details and hyperparameter specifications, refer to the original \ac{BIOT} work~\cite{yang2023biot} or the official implementation at \url{https://github.com/ycq091044/BIOT}.

\subsection{xLSTM Component Configuration}
\label{sec:xlstm_configuration}

\subsubsection{mLSTM Configuration}
\label{sec:mlstm_config}

The \ac{mLSTM} components are configured with the following parameters:

\begin{itemize}
    \item \textbf{Embedding Dimension}: 256
    \item \textbf{Context Length}: 1024 tokens (maximum sequence length)
    \item \textbf{Num Blocks}: 8 (stacked \ac{mLSTM} layers)
    \item \textbf{Dropout}: 0.2 (applied within \ac{mLSTM} cells)
\end{itemize}

The dropout rate of 0.2 provides regularization while maintaining the matrix memory's key-value association capabilities. The \ac{mLSTM} does not include feedforward networks, relying solely on the matrix memory mechanism for feature transformation.

\subsubsection{sLSTM Configuration}
\label{sec:slstm_config}

The \ac{sLSTM} components employ the following configuration:

\begin{itemize}
    \item \textbf{Embedding Dimension}: 256
    \item \textbf{Context Length}: 1024 tokens (maximum sequence length)
    \item \textbf{Num Blocks}: 8 (stacked \ac{sLSTM} layers)
    \item \textbf{Num Heads}: 8 (for multi-head memory operations)
    \item \textbf{Dropout}: 0.2 (applied within \ac{sLSTM} cells)
    \item \textbf{Feedforward}: Enabled with gated activation
    \item \textbf{Projection Factor}: 1.3 (feedforward expansion ratio)
    \item \textbf{Activation Function}: GELU
    \item \textbf{Bias}: False
\end{itemize}

\subsubsection{Combined Architecture Configuration}
\label{sec:combined_config}

The combined \ac{sLSTM}-\ac{mLSTM} architecture employs an eight-block stack with the \textbf{xLSTM[7:1]} configuration:

\begin{itemize}
    \item \textbf{Total Blocks}: 8
    \item \textbf{Block Composition}: 7 \ac{mLSTM} blocks + 1 \ac{sLSTM} block
    \item \textbf{sLSTM Position}: Block index 1 (\ac{mLSTM} → \ac{sLSTM} → \ac{mLSTM} (×6))
    \item \textbf{Dropout}: 0.2 (applied within each cell)
    \item \textbf{Embedding Dimension}: 256 (consistent across all blocks)
    \item \textbf{Context Length}: 1024 tokens (maximum sequence length)
\end{itemize}

This configuration leverages the complementary strengths of both memory mechanisms: \ac{mLSTM}'s parallelizable matrix memory for efficient computation and \ac{sLSTM}'s scalar memory mixing for sequential pattern capture. The single \ac{sLSTM} block positioned early in the stack (block 1) provides crucial memory mixing capabilities while maintaining computational efficiency through predominantly parallel \ac{mLSTM} processing. The 7:1 ratio was identified as the best-performing \ac{mLSTM}-to-\ac{sLSTM} combination through systematic ablation studies in the original xLSTM work~\cite{beck2024xlstm}, where replacing some \ac{sLSTM} blocks with \ac{mLSTM} blocks yielded optimal language modeling performance.

\subsection{Model Parameter Counts}
\label{sec:model_parameters}

The four architecture variants exhibit different parameter counts due to their distinct memory mechanisms and computational requirements:

\begin{table}[h]
\centering
\begin{tabular}{lrr}
\hline
\textbf{Architecture} & \textbf{Parameters} & \textbf{Size (MB)} \\
\hline
Baseline (Linear Transformer) & 3,155,968 & 12.751 \\
mLSTM-only & 3,324,224 & 13.424 \\
sLSTM-only & 2,908,416 & 11.761 \\
Combined (xLSTM[7:1]) & 3,272,248 & 13.216 \\
\hline
\end{tabular}
\caption{Model parameter counts and memory footprint for each architecture variant. All models have comparable complexity, with parameter counts ranging from 2.9M to 3.3M.}
\label{tab:model_parameters}
\end{table}

The \ac{mLSTM}-only configuration has the highest parameter count (3.32M) due to the matrix memory structures in all eight blocks, which require storing key-value matrices. The \ac{sLSTM}-only variant has the fewest parameters (2.91M) as scalar memory mechanisms and gated feedforward networks are more parameter-efficient than matrix operations. The combined xLSTM[7:1] architecture (3.27M parameters) sits between the two pure variants, reflecting its composition of seven \ac{mLSTM} blocks and one \ac{sLSTM} block. The baseline linear Transformer (3.16M parameters) provides a comparable parameter budget for fair comparison.

All models maintain similar complexity levels, ensuring that performance differences reflect architectural advantages rather than capacity disparities. The parameter counts are computed using PyTorch Lightning's model summary functionality, which accounts for all trainable parameters including embeddings, attention mechanisms, memory cells, and classification heads.

\subsection{Classification Head Configuration}
\label{sec:classification_head}

The classification head processes the aggregated sequence representation through mean pooling across all tokens, followed by an \ac{ELU} activation function. The \ac{ELU} activation provides smooth gradients for negative inputs while maintaining identity for positive inputs, facilitating stable optimization. The head outputs a 6-dimensional vector corresponding to the number of event classes in the TUEV dataset. Training employs cross-entropy loss to optimize the model's classification performance.


\section{Training Configuration}
\label{sec:training_configuration}

\subsection{Optimization Parameters}
\label{sec:optimization_parameters}

\textbf{Optimizer Configuration:} The AdamW optimizer combines adaptive learning rates with decoupled weight decay regularization, providing robust optimization for deep neural networks. The learning rate is set to $\alpha = 1 \times 10^{-3}$, with beta parameters $\beta_1 = 0.9$ and $\beta_2 = 0.999$. The epsilon parameter is configured as $\epsilon = 1 \times 10^{-8}$ for numerical stability. A weight decay coefficient of $\lambda = 1 \times 10^{-5}$ provides L2 regularization to prevent overfitting.

\textbf{Batch Configuration:} Training uses a batch size of 128 samples, which balances gradient estimation quality with memory efficiency and training speed. Gradients are applied directly without accumulation. Data loading employs 16 parallel workers with persistent worker mode enabled, meaning worker processes remain alive between epochs to reduce data loading overhead.

\textbf{Training Duration:} Each configuration is trained for a maximum of 100 epochs without early stopping. Instead, the dual checkpointing strategy (described in Section~\ref{sec:checkpointing_config}) captures the best models based on different metrics.

\subsection{Data Loading Configuration}
\label{sec:data_loading}

\textbf{Dataset Split Parameters:} The training data is split at the subject level, with 10\% of training subjects allocated to the validation set. This subject-level splitting ensures that all samples from the same subject remain in the same partition, preventing data leakage and providing a more realistic evaluation of generalization performance. The random seed for splitting is configuration-specific, using values of 10, 42, 123, 456 or 999 depending on the experimental run. Training data is shuffled during loading to prevent batch correlation effects, while validation data maintains its original ordering for consistent evaluation.

\textbf{Data Augmentation:} No data augmentation is applied, as the experiments focus on evaluating base model performance without additional transformations. Signal normalization is performed during preprocessing using per-channel 95th percentile scaling, as described in Section~\ref{sec:preprocessing}.

\subsection{Checkpointing Configuration}
\label{sec:checkpointing_config}

The training process employs a dual checkpointing strategy to capture models that excel according to different criteria. The first checkpoint strategy monitors validation cross-entropy loss, saving the model whenever the loss decreases to a new minimum. The second checkpoint strategy tracks validation balanced accuracy, preserving the model whenever accuracy increases to a new maximum.

Both checkpoints include complete model state, optimizer state, epoch number, and training configuration for reproducibility.

These settings ensure reproducibility across runs with the same seed, though slight variations may occur due to hardware-specific optimizations.


\section{Experimental Execution}
\label{sec:experimental_execution}

\subsection{Parameter Sweep Implementation}
\label{sec:sweep_implementation}

The parameter sweep is implemented through an automated script (\texttt{run\_sweep.py}) that systematically generates and executes all experimental configurations in the design space.

\textbf{Experimental Factors:} The parameter sweep covers three primary factors:
\begin{itemize}
    \item \textbf{Model Architecture:} Four variants (Baseline, sLSTM-only, mLSTM-only, Combined)
    \item \textbf{Temporal Window:} Three configurations (2-2, 3-3, 4-4 seconds before/after events)
    \item \textbf{Random Seeds:} Five seeds (10, 42, 123, 456, 999) for statistical robustness
\end{itemize}

\textbf{Temporal Window Rationale:} The temporal window configurations are designed to explore the impact of contextual information surrounding detected events. Each configuration specifies the number of seconds before and after the central event annotation. For example, a 2-2 configuration creates a 5-second sample: 2 seconds before + 1 second event + 2 seconds after.

The choice of configurations is motivated by both theoretical and practical considerations:
\begin{itemize}
    \item \textbf{Baseline Configuration (2-2):} This 5-second window matches the default configuration used in the original \ac{BIOT} work~\cite{yang2023biot}, where \ac{TUEV} samples are processed with symmetric 2-second contexts around events. Starting from this established baseline enables direct comparison with prior work.
    \item \textbf{Extended Context (3-3):} The 7-second window explores whether additional temporal context improves event characterization by capturing longer-range patterns or anticipatory signals.
    \item \textbf{Maximum Context (4-4):} The 9-second window tests the upper limit of useful temporal context. Further expansion beyond this range was not feasible due to \ac{GPU} memory constraints with the available hardware (RTX 3090 with 24GB).
\end{itemize}

The symmetric temporal windows (secondsBefore = secondsAfter) maintain consistent pre-event and post-event context, simplifying interpretation and enabling fair comparison. This range spans 5-9 seconds total duration, covering clinically plausible contexts while respecting computational constraints.

The multi-seed evaluation addresses the stochastic nature of neural network training. Performance can vary significantly across random initializations due to different local optima, gradient noise, and data shuffling. By evaluating each configuration with five different seeds, the experimental design enables:
\begin{itemize}
    \item Estimation of performance variability
    \item Calculation of confidence intervals
    \item Statistical significance testing between configurations
    \item Identification of robust vs. fragile architectural choices
\end{itemize}

\textbf{Experimental Matrix:} The full factorial design results in $4 \times 5 \times 3 = 60$ training runs, with each run producing two checkpointed models (loss-based and accuracy-based). This comprehensive coverage enables robust statistical analysis and clear identification of optimal configurations for different performance criteria.
%%
%%%%%%%%%%%%%%%%%%%%%%%%%%%%%%%%%%%%%%%%%%%%%%%%%%%%%%%%%%%%%%%%%%%%%%%%%%%%%%%%
